%Input preamble
%Style
\documentclass[12pt]{article}
\usepackage[top=1in, bottom=1in, left=1in, right=1in]{geometry}
\parindent 22pt
\usepackage{fancyhdr}

%Packages
\usepackage{adjustbox}
\usepackage{amsmath}
\usepackage{amsfonts}
\usepackage{amssymb}
\usepackage{bm}
\usepackage[table]{xcolor}
\usepackage{tabu}
\usepackage{color,soul}
\usepackage{makecell}
\usepackage{longtable}
\usepackage{multirow}
\usepackage[normalem]{ulem}
\usepackage{etoolbox}
\usepackage{graphicx}
\usepackage{tabularx}
\usepackage{ragged2e}
\usepackage{booktabs}
\usepackage{caption}
\usepackage{fixltx2e}
\usepackage[para, flushleft]{threeparttablex}
\usepackage[capposition=top,objectset=centering]{floatrow}
\usepackage{subcaption}
\usepackage{pdfpages}
\usepackage{pdflscape}
\usepackage{natbib}
\usepackage{bibunits}
\definecolor{maroon}{HTML}{990012}
\usepackage[colorlinks=true,linkcolor=maroon,citecolor=maroon,urlcolor=maroon,anchorcolor=maroon]{hyperref}
\usepackage{marvosym}
\usepackage{makeidx}
\usepackage{tikz}
\usetikzlibrary{shapes}
\usepackage{setspace}
\usepackage{enumerate}
\usepackage{rotating}
\usepackage{tocloft}
\usepackage{epstopdf}
\usepackage[titletoc]{appendix}
\usepackage{framed}
\usepackage{comment}
\usepackage{xr}
\usepackage{titlesec}
\usepackage{footnote}
\usepackage{longtable}
\newlength{\tablewidth}
\setlength{\tablewidth}{9.3in}
\setcounter{secnumdepth}{4}

\titleformat{\paragraph}
{\normalfont\normalsize\bfseries}{\theparagraph}{1em}{}
\titlespacing*{\paragraph}
{0pt}{3.25ex plus 1ex minus .2ex}{1.5ex plus .2ex}
\makeatletter
\pretocmd\start@align
{%
  \let\everycr\CT@everycr
  \CT@start
}{}{}
\apptocmd{\endalign}{\CT@end}{}{}
\makeatother
%Watermark
\usepackage[printwatermark]{xwatermark}
\usepackage{lipsum}
\definecolor{lightgray}{RGB}{220,220,220}
%\newwatermark[allpages,color=lightgray,angle=45,scale=3,xpos=0,ypos=0]{Preliminary Draft}

%Further subsection level
\usepackage{titlesec}
\setcounter{secnumdepth}{4}
\titleformat{\paragraph}
{\normalfont\normalsize\bfseries}{\theparagraph}{1em}{}
\titlespacing*{\paragraph}
{0pt}{3.25ex plus 1ex minus .2ex}{1.5ex plus .2ex}

\setcounter{secnumdepth}{5}
\titleformat{\subparagraph}
{\normalfont\normalsize\bfseries}{\thesubparagraph}{1em}{}
\titlespacing*{\subparagraph}
{0pt}{3.25ex plus 1ex minus .2ex}{1.5ex plus .2ex}

%Functions
\DeclareMathOperator{\cov}{Cov}
\DeclareMathOperator{\corr}{Corr}
\DeclareMathOperator{\var}{Var}
\DeclareMathOperator{\plim}{plim}
\DeclareMathOperator*{\argmin}{arg\,min}
\DeclareMathOperator*{\argmax}{arg\,max}

%Math Environments
\newtheorem{theorem}{Theorem}
\newtheorem{claim}{Claim}
\newtheorem{condition}{Condition}
\renewcommand\thecondition{C--\arabic{condition}}
\newtheorem{algorithm}{Algorithm}
\newtheorem{assumption}{Assumption}
\renewcommand\theassumption{A--\arabic{assumption}}
\newtheorem{remark}{Remark}
\renewcommand\theremark{R--\arabic{remark}}
\newtheorem{definition}[theorem]{Definition}
\newtheorem{hypothesis}[theorem]{Hypothesis}
\newtheorem{property}[theorem]{Property}
\newtheorem{example}[theorem]{Example}
\newtheorem{result}[theorem]{Result}
\newenvironment{proof}{\textbf{Proof:}}{$\bullet$}

%Commands
\newcommand\independent{\protect\mathpalette{\protect\independenT}{\perp}}
\def\independenT#1#2{\mathrel{\rlap{$#1#2$}\mkern2mu{#1#2}}}
\newcommand{\overbar}[1]{\mkern 1.5mu\overline{\mkern-1.5mu#1\mkern-1.5mu}\mkern 1.5mu}
\newcommand{\equald}{\ensuremath{\overset{d}{=}}}
\captionsetup[table]{skip=10pt}
%\makeindex

\setlength\parindent{20pt}
\setlength{\parskip}{0pt}

\newcolumntype{L}[1]{>{\raggedright\let\newline\\\arraybackslash\hspace{0pt}}m{#1}}
\newcolumntype{C}[1]{>{\centering\let\newline\\\arraybackslash\hspace{0pt}}m{#1}}
\newcolumntype{R}[1]{>{\raggedleft\let\newline\\\arraybackslash\hspace{0pt}}m{#1}}



%Logo
%\AddToShipoutPictureBG{%
%  \AtPageUpperLeft{\raisebox{-\height}{\includegraphics[width=1.5cm]{uchicago.png}}}
%}

\newcolumntype{L}[1]{>{\raggedright\let\newline\\\arraybackslash\hspace{0pt}}m{#1}}
\newcolumntype{C}[1]{>{\centering\let\newline\\\arraybackslash\hspace{0pt}}m{#1}}
\newcolumntype{R}[1]{>{\raggedleft\let\newline\\\arraybackslash\hspace{0pt}}m{#1}}

\newcommand{\mr}{\multirow}
\newcommand{\mc}{\multicolumn}

%\newcommand{\comment}[1]{}


\begin{document}

\doublespacing

\noindent Back-of-the-envelope CBA Calculations \\

\noindent \citet{Kline-Walters_2016_QJE} perform the following exercise:\\

\noindent (i)  Calculate the treatment effect on WPPSI at age 5 for the Head Start Impact Study (HSIS)---HSIS is a one-year-long randomized version of Head Start.\\

\noindent (ii) Monetize this gain using \citet{Chetty_Friedman_etal_2010_HowDoesYour} return to WPPSI at age 5 in terms of net present value of earnings at age-27. \citet{Chetty_Friedman_etal_2010_HowDoesYour} calculations indicate that a 1 standard deviation gain in WPPSI at age 5 implies a $13.1\%$ increase in the net present value of earnings at age 27.\\

\noindent (iii) Calculate the benefit-to-cost ratio by dividing the gain in ``(ii)" by the total costs of the program. This calculation is based on giving a value to net-present value of earnings at age 27 of $\$385,907.17$ to the control-group participants, which is provided by \citet{Chetty_Friedman_etal_2010_HowDoesYour}. The treatment-group participants are given $13.1\%$ increase on this amount. \citet{Kline-Walters_2016_QJE} calculation indicates that the benefit-to-cost ratio of HSIS ranges between $1.50$ and $1.84$.\\ 

\noindent We perform three exercises related to this.\\

\noindent First, we perform the same calculation described before, but use the gain in WPPSI at age 5 and the program costs of ABC/CARE. This gain amounts to .48 of a standard deviation. Using the same return to WPPSI in terms of earnings and the same net present value of income as in\citet{Chetty_Friedman_etal_2010_HowDoesYour} we obtain a benefit-to-cost of $.45$.\\

\noindent Second, we calculate the return of WPPSI in terms of net present value of earnings. As opposed to  \citet{Chetty_Friedman_etal_2010_HowDoesYour}, we consider the net present value of life-cycle earnings and not the net-present value of earnings at age 27. We provide details on how to obtain this life-cycle value in the main body of the paper. A 1 standard deviation gain in WPPSI at age 5 implies a $84.6\%$ increase in life-cycle net present value of earnings. The net present value of life-cycle earnings for the control group is $\$298,611.90$ (in 2014 USD). Applying the ``return to WPPSI'' in terms of life-cycle earnings to this quantity to the participants of the treatment group, we are able to produce a benefit-to-cost ratio. The calculation indicates that this number is: $1.48$.\\

\noindent Finally, we calculate the return of WPPSI in terms of life-cycle net present value of all the components we monetize in the paper. A 1 standard deviation gain in WPPSI at age 5 implies an increase in the life-cycle net present value of all components of $119.03\%$. The life-cycle net present value of all components for the control group is $\$753,692.98$. Applying the ``return to WPPSI'' in terms of the life-cycle earnings to this quantity to the participants of the treatment group yields a benefit-to-cost ratio of $5.15$.\footnote{The difference between this number and the number we report in the main paper due to leaving out CARE from this exercise for simplicity.}\\

%References
\singlespace
\bibliographystyle{chicago}
\bibliography{heckman}
\end{document}